\documentclass[11pt,a4paper]{article}
\usepackage[utf8]{inputenc}
\usepackage{amsmath,amssymb,amsthm}
\usepackage{listings}
\usepackage{xcolor}
\usepackage[margin=1in]{geometry}
\usepackage{hyperref}

\newtheorem{theorem}{Theorem}

\lstset{
  language=C++,
  basicstyle=\ttfamily\small,
  keywordstyle=\color{blue}\bfseries,
  commentstyle=\color{green!60!black},
  stringstyle=\color{red},
  numbers=left,
  numberstyle=\tiny\color{gray},
  breaklines=true,
  frame=single,
  tabsize=4,
  showstringspaces=false
}

\title{IOI 2004 -- Polygon Triangulation}
\author{Solution}
\date{}

\begin{document}
\maketitle

\section{Problem Statement Summary}
Given a convex polygon with $N$ vertices, each having a weight $w_i$,
triangulate it (divide into $N - 2$ triangles using non-crossing
diagonals) to maximize the total score. The score of triangle
$(i, k, j)$ is $w_i \cdot w_k \cdot w_j$.

\section{Solution: Interval DP}

This is structurally identical to optimal matrix-chain multiplication.

\subsection{Recurrence}
Fix edge $(i, j)$. Any triangulation of the sub-polygon
$i, i{+}1, \ldots, j$ includes exactly one triangle containing edge
$(i, j)$, with some third vertex $k$ ($i < k < j$). This gives:
\[
  \mathrm{dp}[i][j] = \max_{i < k < j}
  \bigl(\mathrm{dp}[i][k] + \mathrm{dp}[k][j] + w_i \cdot w_k \cdot w_j\bigr).
\]
Base case: $\mathrm{dp}[i][i{+}1] = 0$ (an edge, no triangle).
Answer: $\mathrm{dp}[0][N{-}1]$.

\begin{theorem}
The interval DP correctly computes the optimal triangulation because
every triangulation of a convex polygon is uniquely determined by the
choice of third vertex for each edge's triangle, and the subproblems
are independent.
\end{theorem}

\section{C++ Implementation}

\begin{lstlisting}
#include <bits/stdc++.h>
using namespace std;

int main() {
    ios::sync_with_stdio(false);
    cin.tie(nullptr);

    int N;
    cin >> N;

    vector<long long> w(N);
    for (int i = 0; i < N; i++) cin >> w[i];

    // dp[i][j] = max score of triangulating sub-polygon i..j
    vector<vector<long long>> dp(N, vector<long long>(N, 0));

    for (int len = 3; len <= N; len++) {
        for (int i = 0; i + len - 1 < N; i++) {
            int j = i + len - 1;
            dp[i][j] = LLONG_MIN;
            for (int k = i + 1; k < j; k++) {
                long long val = dp[i][k] + dp[k][j] + w[i] * w[k] * w[j];
                dp[i][j] = max(dp[i][j], val);
            }
        }
    }

    cout << dp[0][N - 1] << "\n";
    return 0;
}
\end{lstlisting}

\section{Complexity Analysis}
\begin{itemize}
  \item \textbf{Time:} $O(N^3)$. Three nested loops: interval length,
        start position, and split point.
  \item \textbf{Space:} $O(N^2)$.
\end{itemize}

For $N \le 500$ this runs comfortably. If the polygon is given as a
ring (any edge can be the base), ``break the ring'' by fixing one edge
or by considering all $N$ rotations.

\end{document}
